% Appendix C: Per-experiment Measurements

\section{Per-Experiment Measurements}
\label{app:measurements}

This appendix collects the regtest measurements underlying
Section~\ref{sec:empirical} at greater granularity than the main
body's summary tables permit. Each value is structurally
deterministic (range~=~0 across repeated runs on the reference
implementations: \texttt{simple-ctv-vault}, \texttt{pymatt},
\texttt{opvault-demo}, our \texttt{cat-csfs-vault}, and our
\texttt{simple-simplicity-vault}).

\subsection{Revault Amplification}
\label{app:meas:revault}

Per-step vsize for nine successive partial withdrawals of
$5{,}000{,}000$~sats each from an initial $49{,}999{,}900$-sat
vault. Per-step size is constant across all nine steps for every
design, confirming structural determinism.

\begin{table}[h]
\centering
\caption{Per-step vsize for repeated partial unvaults with
revault. \opccv{} and \opvault{} measure their native
\texttt{trigger\_and\_revault}; \opctv{}, \catcsfs{}, and Simplicity do
not support native revault---the figures listed show the
single-withdrawal cost that repeats per step in a full-cycle
emulation (unvault $\rightarrow$ withdraw $\rightarrow$
re-deposit).}
\label{tab:revault_appendix}
\small
\begin{tabular}{@{}lrrl@{}}
\toprule
Design   & Per-step vB & Cumulative @ 9 steps (vB) & Revault model \\
\midrule
\opccv{}       & 111 & 999    & Native partial unvault \\
\opvault{} & 121 & 1{,}089 & Native partial unvault \\
\opctv{}       & 152 & 1{,}368 & Full-cycle emulation \\
\catcsfs{}  & 210 & 1{,}890 & Full-cycle emulation \\
Simplicity & 386 & 3{,}474 & Full-cycle emulation (Elements serialisation) \\
\bottomrule
\end{tabular}
\end{table}

\subsection{Batching Sweep}
\label{app:meas:batching}

Total transaction vsize as a function of batch size $N$ for a
single trigger aggregating $N$ vault UTXOs. \opccv{} and \opvault{}
amortise fixed overhead across inputs (sub-linear scaling); \opctv{},
\catcsfs{}, and Simplicity lack a spec-level batching primitive, so
the natural cost is $N$ independent transactions and the totals
grow strictly linearly.

\begin{table}[h]
\centering
\caption{Batched trigger vsize measurements at batch sizes
$N \in \{1, 2, 3, 5, 10, 20, 50\}$. Regression on the \opccv{} and
\opvault{} rows yields $\text{vsize}_{\text{trigger}}(N) =
\alpha_t + \beta_t \cdot N$ with
$(\alpha_t, \beta_t)_{\mathrm{OP\_CCV}} \approx (53, 100)$\vB{} and
$(\alpha_t, \beta_t)_{\opvault{}} \approx (198, 94)$\vB{}. The
recovery-side linear fits used in
Section~\ref{sec:empirical:batching} are
$(\alpha, \beta)_{\mathrm{OP\_CCV}} = (54, 68)$\vB{} and
$(\alpha, \beta)_{\opvault{}} = (154, 92)$\vB{}.}
\label{tab:batching_appendix}
\footnotesize
\setlength{\tabcolsep}{4pt}
\begin{tabular}{@{}lrrrrrrr@{}}
\toprule
Design    & $N{=}1$ & $N{=}2$ & $N{=}3$ & $N{=}5$ & $N{=}10$ & $N{=}20$ & $N{=}50$ \\
\midrule
\opccv{}       & 154 & 254  & 355   & 555    & 1{,}056  & 2{,}059  & 5{,}066 \\
\opvault{} & 292 & 385  & 479   & 667    & 1{,}135  & 2{,}073  & 4{,}885 \\
\midrule
\opctv{}       & 94  & 188  & 282   & 470    & 940      & 1{,}880  & 4{,}700 \\
\catcsfs{}  & 221 & 442  & 663   & 1{,}105 & 2{,}210 & 4{,}420  & 11{,}050 \\
Simplicity & 298 & 596 & 894   & 1{,}490 & 2{,}980 & 5{,}960  & 14{,}900 \\
\bottomrule
\end{tabular}
\end{table}

\subsection{Fee-Rate Sensitivity Projections}
\label{app:meas:fees}

Projected on-chain lifecycle cost
$C_{\text{lifecycle}}(r) = \text{vsize}_{\text{total}} \times r$
for each Bitcoin design at six historically observed fee rates.
Fee-pinning cost is $r$-invariant. Watchtower-exhaustion
feasibility is fee-dependent: the rightmost column gives the
number of per-UTXO splits an attacker needs to drive a $0.5$~BTC
\opccv{} vault below dust. The ranking of lifecycle totals is
fee-invariant because vsize ratios are constant.

\begin{table}[h]
\centering
\caption{Projected lifecycle cost (sats) across Bitcoin designs
and the \opccv{} split-to-dust count, at $V = 0.5$~BTC. Simplicity
runs on Elements and uses an incomparable fee mechanism (explicit
fee outputs), so it is excluded from Bitcoin-fee projections.}
\label{tab:fee_sensitivity_appendix}
\footnotesize
\begin{tabular}{@{}rrrrrrr@{}}
\toprule
& \multicolumn{4}{c}{Lifecycle cost (sats)} & Fee-pin & \opccv{} splits \\
\cmidrule(lr){2-5}
$r$ (\satvB{}) & \opctv{} & \opccv{} & \opvault{} & \catcsfs{} & (all $r$) & to dust \\
\midrule
1   &     368 &     565 &     567 &     553 & 14{,}300 & 409{,}836 \\
10  &   3{,}680 &   5{,}650 &   5{,}670 &   5{,}530 & 14{,}300 &  40{,}984 \\
50  &  18{,}400 &  28{,}250 &  28{,}350 &  27{,}650 & 14{,}300 &   8{,}197 \\
100 &  36{,}800 &  56{,}500 &  56{,}700 &  55{,}300 & 14{,}300 &   4{,}098 \\
300 & 110{,}400 & 169{,}500 & 170{,}100 & 165{,}900 & 14{,}300 &   1{,}366 \\
500 & 184{,}000 & 282{,}500 & 283{,}500 & 276{,}500 & 14{,}300 &     820 \\
\bottomrule
\end{tabular}
\end{table}

\subsection{Address Reuse}
\label{app:meas:reuse}

A second deposit to the same \opctv{} vault address creates a UTXO
whose amount does not match the committed template hash; the
second UTXO is permanently unspendable via the vault path
(confirmed on regtest). \opccv{}, \opvault{}, \catcsfs{}, and Simplicity
all handle repeated deposits safely: \opccv{} and \opvault{} validate
each UTXO per-instance against the contract; \catcsfs{} generates a
unique vault plan per deposit (distinct P2TR addresses);
Simplicity creates a fresh plan per instance. This is a
design-level property of \opctv{}'s input-count-and-amount commitment,
not an implementation artefact.

\subsection{Measurements Not Captured}
\label{app:meas:limits}

Regtest mines blocks on demand and has no fee market, so temporal
race dynamics (TM3 attacker-vs-watchtower timing, fee-pinning
propagation delays) cannot be measured on the harness. Structural
vsize is identical on regtest and mainnet for any fixed
transaction structure, so the lifecycle, batching, and revault
tables above transfer directly; the $r^\dagger$ thresholds of
Section~\ref{sec:empirical:batching} are derived by projection,
not timing.

\providecommand{\todo}[2][]{\textbf{[TODO: #2]}}
